%%%%%%%%%%%%%%%%%%%%%%%%%%%%%%%%%%%%%%%%%%%%%%%%%%%%%%%%%%%%%%%%%%%%%%%%
%                                                                      %
%     File: Thesis_Conclusions.tex                                     %
%     Tex Master: Thesis.tex                                           %
%                                                                      %
%     Author: Francisco Castro                                         %
%     Last modified :  24 Fev 2026                                     %
%                                                                      %
%%%%%%%%%%%%%%%%%%%%%%%%%%%%%%%%%%%%%%%%%%%%%%%%%%%%%%%%%%%%%%%%%%%%%%%%

\chapter{Conclusions}
\label{chapter:conclusions}

The work that was developped in this thesis has met the majority of the goals that were stated in section \ref{section:objectives}. Many traversability maps were successfully generated by capturing the adaptability of the mobile robot to the irregular terrain. Then, optimal paths were correctly computed by avoiding the high cost positions of the respective traversability map. The type of traversability map and respective potential field depend on the hull complexity and the parameters  of the traversability map computation algorithm.

Secondly, extensive testing of the compliance of the \acrshort{nmpc} algorithms to real-time requirements was carried out, which confidently ensured this algorithm is suitable for real-time applications in rough terrain. The mobile robot has rarely been on potentially unsafe poses or scenarios that could compromise either the controller feasibility or the vehicle integrity. The controller was refined to prevent slippage as much as possible. Moreover, safe areas were assigned to each node of the control horizon  so that the mobile robot is free from any collision at any time.

One of the main deficiencies of this thesis is that the presented \acrshort{nmpc} algorithms were not taken to a simulation environment or a real-world experiment. Apart from that, an effective traction control simulation real-time bed was not used, due to the fact that it was not possible to successfully capture the dynamics regime of the \acrshort{ode} physics engine of Gazebo. Another goal that was meant to be achieved and was left behind was the integration of the global, safe, optimal path generated in chapter \ref{chapter:globalpathplanner} with the \acrshort{nmpc} algorithm. 

The virtual speed control module has a lot of room for improvement. In this thesis, the virtual speed control module has only considered the path curvature and distance to the first reference node, but other parameters such as the safe balls radius can be included.

Even though the authour showed that the controller was able keep the mobile robot away from collisions, it is assumed that the provided deformed reference shall not be within unsafe space. One possible work-around for this is to increase the repulsive fields magnitude, but the problem of local minimums remains to solve.

% ----------------------------------------------------------------------
\section{Future Work}
\label{section:future}

From this short review of the dissertation main technical deliverables and deficits, this section provides high-level directions to future research.

This thesis has just presented controllers that include the dynamics model of rigid vehicles. However, irregular terrains are typically crossed by articulated mobile robots\footnote{Leo-rover: https://www.leorover.tech/ (accessed on 23 of March 2026)}.

If only one source is selected at a time, the global path planner needs to recompute the Eikonal equation \textit{in-situ} in order to provide a path to the new goal. If the Eikonal equation can work with several sources at once, only one computation of the potential field is required. In the literature, this approach is called \textit{Fast marching firework method}.

Gazebo has not proven to be a reliable wheeled robot simulator test-bed for rough terrains. Project Chrono\footnote{Project Chrono: https://projectchrono.org/ (accessed on 23 of March 2026)} (open-source) and Vortex Studio\footnote{Vortex Studio: https://www.cm-labs.com/en/vortex-studio/ (accessed on 23 of March 2026)} (commercial software) are more matured simulators for these terrains. The interaction between rigid wheels and deformable terrain are modeled by them. Reportedly, \acrshort{ros} bridges to these softwares exist.

Moreover, a variable, controlled virtual speed between consecutive reference nodes can be implemented as well.

Currently, the safe ball radius is not optimal. In the top right picture of figure \ref{fig:safeballsexamplessolver1}, one can see that there is still room to increase the safe balls area. Therefore, more refined search methods could be envisioned to tackle this deficiency. Moreover, only circled obstacles were considered. Then, a reformulation of algorithms \nameref{alg:computefreespace} and \nameref{alg:gendeformedref} that copes with more complex obstacles is a possible line of research. A more robust generation method of references \textit{in-situ} shall be envisioned as the artificial potential field approach may fall on local minimums.

The \textit{online} determination of the terrain parameters is a typical research subject that was not addressed in this thesis. Optimal parameter estimation can be used to estimate them \cite{diehlnumericalocp}. Then, the determined parameters can be fed to the \acrshort{nmpc} module, and the vehicle's response to a controller that is unaware of the exact terrain parameters can be studied.